\documentclass[12pt,a4paper]{article}
\usepackage[T1]{fontenc}
\usepackage[utf8]{inputenc}
\usepackage[spanish]{babel}
\usepackage{geometry}
\usepackage{booktabs}
\usepackage{xcolor}
\usepackage{listings}
\usepackage{hyperref}
\usepackage{enumitem}
\usepackage{titlesec}
\usepackage{parskip}

\geometry{margin=2.5cm}

\hypersetup{
  colorlinks=true,
  linkcolor=black,
  urlcolor=blue!70!black,
}

\lstdefinestyle{bash}{
  language=bash,
  basicstyle=\small\ttfamily,
  backgroundcolor=\color{gray!10},
  frame=single,
  rulecolor=\color{gray!40},
  breaklines=true,
  showstringspaces=false,
}

\title{%
  \textbf{Micrófono interno sin funcionar}\\[0.4em]
  \large Diagnóstico y solución en el ASUS ROG Zephyrus G14 GA403UV\\
  con CachyOS + Omarchy (PipeWire / WirePlumber)
}
\author{pc-config}
\date{Septiembre de 2026}

\begin{document}

\maketitle
\tableofcontents
\newpage

% ─────────────────────────────────────────────
\section{Resumen}

El micrófono interno del \textbf{ASUS ROG Zephyrus G14 GA403UV} aparecía
identificado por el sistema (visible en \texttt{arecord}, \texttt{wpctl} y
\texttt{pavucontrol}) pero \textbf{no capturaba audio}: OBS Studio no detectaba
sonido y en \texttt{alsamixer} no se veían niveles.

La causa no era el firmware, ni PipeWire, ni el coprocesador de audio AMD
(ACP/SOF): era el \textbf{mux de entrada del codec Realtek ALC285}, que estaba
seleccionando el micrófono del \emph{jack} externo (sin nada conectado) en
lugar de los micrófonos internos.

La solución son tres comandos de \texttt{amixer} más el guardado del estado
ALSA para que sobreviva reinicios.

% ─────────────────────────────────────────────
\section{Entorno del sistema}

\begin{center}
\begin{tabular}{ll}
  \toprule
  Componente & Versión / valor \\
  \midrule
  Portátil            & ASUS ROG Zephyrus G14 GA403UV \\
  Sistema operativo   & CachyOS (base Arch Linux) \\
  Entorno de escritorio & Omarchy (Hyprland + Wayland) \\
  Kernel activo       & \texttt{linux-cachyos} 7.0.9-1 \\
  Codec de audio      & Realtek ALC285 (tarjeta 2, HDA) \\
  Amplificadores      & 2 × Cirrus Logic CS35L56 (I2C) para altavoces \\
  Servidor de audio   & PipeWire 1.6.8 + WirePlumber 0.5.17 \\
  Firmware SOF        & \texttt{sof-firmware} 2025.12.2-1 \\
  UCM                 & \texttt{alsa-ucm-conf} 1.2.16.1-1 \\
  Utilidades ALSA     & \texttt{alsa-utils} 1.2.16-1 \\
  \bottomrule
\end{tabular}
\end{center}

El equipo tiene \textbf{cuatro} dispositivos de audio PCI:

\begin{lstlisting}[style=bash]
$ lspci -nn | grep -i audio
01:00.1 Audio device: NVIDIA Corporation AD107 HD Audio Controller [10de:22be]
65:00.1 Audio device: AMD/ATI Radeon High Definition Audio [1002:1640]
65:00.5 Multimedia controller: AMD Audio Coprocessor [1022:15e2] (rev 63)
65:00.6 Audio device: AMD Ryzen HD Audio Controller [1022:15e3]
\end{lstlisting}

El micrófono interno \textbf{cuelga del codec ALC285} (función HDA en
\texttt{65:00.6}), no del ACP. Por eso no hace falta tocar SOF ni
\texttt{snd\_sof\_amd\_*} para este problema.

% ─────────────────────────────────────────────
\section{Síntoma}

\begin{itemize}
  \item El micrófono figura como dispositivo de captura válido
        (\texttt{arecord -l} muestra \texttt{card 2: Generic\_1 [HD-Audio
        Generic], device 0: ALC285 Analog}).
  \item La fuente existe en PipeWire:
        \texttt{alsa\_input.pci-0000\_65\_00.6.analog-stereo}.
  \item Al grabar o al usar OBS, \textbf{silencio absoluto} (RMS 0).
  \item En \texttt{alsamixer} no se observan niveles de captura.
\end{itemize}

% ─────────────────────────────────────────────
\section{Diagnóstico}

\subsection{Comprobar el stack de audio}

\begin{lstlisting}[style=bash]
# Dispositivos de captura
arecord -l

# Tarjetas ALSA
cat /proc/asound/cards

# Estado de PipeWire (fuentes y sinks)
wpctl status

# Puertos del card del codec
pactl list cards | grep -E 'alsa_card|analog-input|Ports:'
\end{lstlisting}

Salida relevante observada:

\begin{lstlisting}[style=bash]
**** List of CAPTURE Hardware Devices ****
card 2: Generic_1 [HD-Audio Generic], device 0: ALC285 Analog [ALC285 Analog]
  Subdevices: 1/1

# Puertos del card 65:00.6
analog-input-internal-mic: Microfono interno (priority: 8900, unknown)
analog-input-mic: Microfono (priority: 8700, not available)
\end{lstlisting}

\subsection{El control real: \texttt{Capture Source}}

En este codec el mux de entrada no se expone como un simple switch de
alsa\-mixer, sino como un \textbf{control enumerado}:

\begin{lstlisting}[style=bash]
$ amixer -c 2 cget numid=19
numid=19,iface=MIXER,name='Capture Source'
  ; type=ENUMERATED,access=rw------,values=1,items=4
  ; Item #0 'Internal Mic'
  ; Item #1 'Internal Mic 1'
  ; Item #2 'Mic'
  ; Item #3 'Mic 1'
  : values=2        # <-- 'Mic' = jack externo (SIN conectar)
\end{lstlisting}

Los micrófonos internos son dos pines digitales del codec:

\begin{lstlisting}[style=bash]
$ grep -E 'Node 0x12|Node 0x13' -A2 /proc/asound/card2/codec#0
Node 0x12 [Pin Complex] wcaps 0x40040b: Stereo Amp-In
  Pin Default 0x90a60140: [Fixed] Mic at Int N/A   # mic interno 1
Node 0x13 [Pin Complex] wcaps 0x40040b: Stereo Amp-In
  Pin Default 0x90a60150: [Fixed] Mic at Int N/A   # mic interno 2
\end{lstlisting}

Además, la ganancia estaba al máximo: \texttt{Internal Mic Boost = 30\,dB}
(3 de 3), lo que saturaba cualquier señal.

% ─────────────────────────────────────────────
\section{Causa raíz}

\begin{enumerate}
  \item \texttt{Capture Source} apuntaba al ítem \texttt{Mic}
        (micrófono del \emph{jack} combo, nodo externo), que no tiene nada
        conectado $\Rightarrow$ captura de silencio.
  \item Aun seleccionando la entrada correcta, el \emph{boost} de hardware a
        30\,dB produce \emph{clipping} masivo (se midió 27\,\% de muestras
        recortadas).
\end{enumerate}

No es un fallo de firmware, ni de PipeWire/WirePlumber, ni del ACP/SOF.
El estado guardado en \texttt{/var/lib/alsa/asound.state} simplemente tenía
seleccionada la entrada equivocada.

% ─────────────────────────────────────────────
\section{Solución}

\subsection{Comandos}

\begin{lstlisting}[style=bash]
# 0) Confirmar la tarjeta del codec (aqui: card 2, ALC285)
cat /proc/asound/cards

# 1) Seleccionar el microfono interno (nodos 0x12/0x13)
amixer -c 2 cset iface=MIXER,name='Capture Source' 'Internal Mic'

# 2) Ganancia limpia: boost 0 dB y captura al 84 %
amixer -c 2 sset 'Internal Mic Boost',0 0
amixer -c 2 sset 'Capture' 84%

# 3) Guardar el estado para que sobreviva reinicios
sudo alsactl store
\end{lstlisting}

\begin{itemize}
  \item \texttt{alsactl store} se puede lanzar sin terminal root con
        \texttt{pkexec alsactl store}; en Omarchy salta un diálogo gráfico de
        polkit.
  \item Si la tarjeta no fuese la 2, usar el nombre:
        \texttt{amixer -c Generic\_1 ...} o ajustar el número según
        \texttt{/proc/asound/cards}.
\end{itemize}

\subsection{Nota sobre la sintaxis de amixer}

Con \texttt{alsa-utils} 1.2.16, para los switches exclusivos del codec
(\texttt{Internal Mic}, \texttt{Mic}, etc.) la forma \texttt{sset 'Internal
Mic' on} falla con \texttt{amixer: Invalid command!}. Hay que usar el enum
(\texttt{cset ... 'Capture Source' 'Internal Mic'}) o \texttt{toggle}.
El comando documentado arriba con \texttt{cset} es el correcto y verificado.

\subsection{Persistencia}

El servicio \texttt{alsa-restore.service} (estático, se ejecuta en cada
arranque) restaura los controles desde:

\begin{lstlisting}[style=bash]
/var/lib/alsa/asound.state
\end{lstlisting}

Verificar que quedó bien guardado:

\begin{lstlisting}[style=bash]
grep -A10 "name 'Capture Source'" /var/lib/alsa/asound.state
# Debe contener:  value 'Internal Mic'
\end{lstlisting}

\subsection{Solución permanente: servicio de usuario}

Además del estado ALSA, se instaló un servicio de usuario que reafirma el
micrófono interno en cada inicio de sesión (por si una actualización de kernel
o de PipeWire revierte el estado). Ya está habilitado:

\begin{lstlisting}[style=bash]
systemctl --user status internal-mic.service
\end{lstlisting}

\texttt{\textasciitilde/.local/bin/internal-mic-fix}:

\begin{lstlisting}[style=bash]
#!/usr/bin/env bash
# Fija los microfonos internos del ALC285, esperando al codec.
set -u
AMIXER=/usr/bin/amixer
for _ in $(seq 1 30); do
    for card in 0 1 2 3 4 5; do
        if "$AMIXER" -c "$card" cget iface=MIXER,name='Capture Source' \
             >/dev/null 2>&1; then
            "$AMIXER" -c "$card" cset iface=MIXER,name='Capture Source' \
                 'Internal Mic' >/dev/null 2>&1
            "$AMIXER" -c "$card" sset 'Internal Mic Boost',0 0 \
                 >/dev/null 2>&1
            exit 0
        fi
    done
    sleep 1
done
exit 1
\end{lstlisting}

\texttt{\textasciitilde/.config/systemd/user/internal-mic.service}:

\begin{lstlisting}[style=bash]
[Unit]
Description=Fijar microfono interno (ALC285) como fuente de captura
After=pipewire.service wireplumber.service

[Service]
Type=oneshot
ExecStart=%h/.local/bin/internal-mic-fix
RemainAfterExit=yes

[Install]
WantedBy=default.target
\end{lstlisting}

\begin{lstlisting}[style=bash]
chmod +x ~/.local/bin/internal-mic-fix
systemctl --user daemon-reload
systemctl --user enable --now internal-mic.service
\end{lstlisting}

% ─────────────────────────────────────────────
\section{Verificación}

Grabar unos segundos y medir la señal (hablar durante la grabación):

\begin{lstlisting}[style=bash]
timeout 8 pw-record \
  --target alsa_input.pci-0000_65_00.6.analog-stereo \
  /tmp/test-mic.wav
\end{lstlisting}

Análisis rápido del WAV:

\begin{lstlisting}[style=bash]
python3 -c "
import wave, struct
w = wave.open('/tmp/test-mic.wav','rb')
data = w.readframes(w.getnframes())
s = struct.unpack('<%dh' % (len(data)//2), data)
peak = max(abs(x) for x in s)
rms = (sum(x*x for x in s)/len(s))**0.5
clip = sum(1 for x in s if abs(x) >= 32000)
print(f'peak={peak} ({peak/32768*100:.1f}%) rms={rms:.0f} '
      f'clip={clip/len(s)*100:.2f}%')
"
\end{lstlisting}

Resultado esperado con voz normal:

\begin{itemize}
  \item \texttt{peak} entre 5\,\% y 50\,\% (sin llegar a 100\,\%).
  \item \texttt{rms} mayor a 0.
  \item \texttt{clip} cercano a 0\,\%.
\end{itemize}

Si aparece \texttt{peak=32768} y \texttt{clip} alto, bajar la ganancia
(\texttt{Capture} o el boost). Si la señal es muy baja, subirla primero en la
aplicación (OBS permite filtros de ganancia) antes que el boost de hardware.

% ─────────────────────────────────────────────
\section{Uso en aplicaciones (OBS)}

\begin{enumerate}
  \item Añadir fuente \textbf{Audio Input Capture} (\emph{Captura de entrada
        de audio}).
  \item Seleccionar el dispositivo:
        \texttt{Ryzen HD Audio Controller Estéreo analógico}
        (internamente \texttt{alsa\_input.pci-0000\_65\_00.6.analog-stereo}).
  \item En el mezclador, activar el monitoreo o verificar el vúmetro al
        hablar.
\end{enumerate}

\subsection{Advertencia: la fuente puede apuntar al headset Bluetooth}

Caso real encontrado tras un reinicio: el sistema estaba perfecto (el micrófono
interno capturaba bien), pero OBS seguía sin sonido porque su fuente
\texttt{Mic/Aux} (\texttt{AuxAudioDevice1}) había quedado apuntando al
micrófono de un headset Bluetooth:

\begin{lstlisting}[style=bash]
"device_id": "bluez_input.74:B7:E6:37:11:11"
\end{lstlisting}

Si el headset está apagado o desconectado, esa fuente no existe y OBS captura
silencio. Verificar en \emph{Propiedades} de la fuente que el dispositivo sea
el micrófono interno (o \texttt{Predeterminado}). La configuración se guarda
en:

\begin{lstlisting}[style=bash]
~/.config/obs-studio/basic/scenes/<escena>.json   # AuxAudioDevice1
\end{lstlisting}

Cerrar OBS antes de editar el JSON a mano, y hacer copia de seguridad.

% ─────────────────────────────────────────────
\section{Notas y mantenimiento}

\subsection{Si se conecta un headset con micrófono}

El ítem \texttt{Mic} del enum corresponde al micrófono del \emph{jack}. Para
usarlo:

\begin{lstlisting}[style=bash]
amixer -c 2 cset iface=MIXER,name='Capture Source' 'Mic'
# ... y al desenchufar, volver a 'Internal Mic'
\end{lstlisting}

También se puede cambiar desde la interfaz gráfica de PipeWire (pavucontrol /
\texttt{wpctl}), seleccionando el puerto \emph{Micrófono} en la pestaña de
dispositivos de entrada.

\subsection{Por qué no hay que tocar SOF/ACP}

Aunque los módulos \texttt{snd\_sof\_amd\_acp63} y \texttt{snd\_pci\_ps} estén
cargados, el ACP (\texttt{65:00.5}) no expone ninguna tarjeta de captura en
este equipo y \textbf{no es necesario}: los micrófonos internos están en los
pines 0x12/0x13 del ALC285. No instalar ni forzar controladores adicionales
por este problema.

\subsection{Diagnóstico rápido de 30 segundos}

\begin{lstlisting}[style=bash]
# 1) Estado del mux (debe decir Internal Mic)
amixer -c 2 cget numid=19

# 2) Niveles
amixer -c 2 sget 'Internal Mic Boost',0 | grep 'Front Left'
amixer -c 2 sget 'Capture' | grep 'Front Left'

# 3) Fuente visible en PipeWire
wpctl status | grep -A3 'Sources:'
\end{lstlisting}

\subsection{Referencia de controles}

\begin{center}
\begin{tabular}{lll}
  \toprule
  Control & Rango & Valor recomendado \\
  \midrule
  \texttt{Capture Source} (numid 19) & enum 4 ítems & \texttt{Internal Mic} \\
  \texttt{Internal Mic Boost},0 & 0--3 (0--30\,dB) & 0 \\
  \texttt{Capture} & 0--63 & 84\,\% (53) \\
  \texttt{Master} & 0--87 & 100\,\% (87) \\
  \bottomrule
\end{tabular}
\end{center}

\end{document}
